% ============================================================
%  BAB 6 – PENUTUP
%  Compile standalone:  pdflatex bab-6.tex
%  Compile as part of thesis: pdflatex main.tex
% ============================================================
\documentclass[../../thesis]{subfiles}
\usepackage{hyphenat}
\usepackage{iftex}

% ── Encoding, Language & Fonts ───────────────────────────────
\ifXeTeX
    \usepackage{polyglossia}
    \setmainlanguage{bahasai}
    \usepackage{fontspec}
    \setmainfont{TeX Gyre Termes}
    \setsansfont{TeX Gyre Heros}
    \setmonofont[Scale=0.88]{TeX Gyre Cursor}
\else
    \usepackage[utf8]{inputenc}
    \usepackage[T1]{fontenc}
    \usepackage[indonesian]{babel}
    \usepackage{mathptmx}
    \usepackage{helvet}
    \usepackage{courier}
\fi

% ── Page geometry (A4, UI thesis margins) ────────────────────
\usepackage[
  a4paper,
  top    = 2.54cm,
  bottom = 2.54cm,
  left   = 3.17cm,
  right  = 2.54cm
]{geometry}

% ── Line spacing & paragraph ─────────────────────────────────
\usepackage{setspace}
\onehalfspacing
\setlength{\parindent}{1.27cm}
\setlength{\parskip}{0pt}

% ── Microtype & justification ────────────────────────────────
\usepackage{microtype}
\sloppy

% ── Headings ─────────────────────────────────────────────────
\usepackage{titlesec}

\titleformat{\chapter}[block]
  {\centering\rmfamily\bfseries\large}{}
  {0pt}{\MakeUppercase}
\titlespacing*{\chapter}{0pt}{0pt}{1.5em}

\titleformat{\section}
  {\rmfamily\bfseries\normalsize}{\thesection}{1em}{}
\titlespacing*{\section}{0pt}{1.2em}{0.6em}

\titleformat{\subsection}
  {\rmfamily\bfseries\normalsize}{\thesubsection}{1em}{}
\titlespacing*{\subsection}{0pt}{1em}{0.5em}

\titleformat{\subsubsection}
  {\rmfamily\bfseries\normalsize}{\thesubsubsection}{1em}{}
\titlespacing*{\subsubsection}{0pt}{0.8em}{0.4em}

% ── Math ─────────────────────────────────────────────────────
\usepackage{amsmath, amssymb}

% ── Lists ────────────────────────────────────────────────────
\usepackage{enumitem}
\setlist[enumerate,1]{
  label      = \arabic*.,
  leftmargin = 1.27cm,
  itemsep    = 4pt,
  parsep     = 0pt,
  topsep     = 4pt
}
\setlist[enumerate,2]{
  label      = \alph*.,
  leftmargin = 1.27cm,
  itemsep    = 3pt,
  parsep     = 0pt,
  topsep     = 2pt
}

% ── Hyperlinks ───────────────────────────────────────────────
\usepackage[hidelinks]{hyperref}

% Bibliography setup is inherited from main.tex (biblatex + references.bib)
% when this file is compiled standalone via the subfiles class.

% ============================================================
\begin{document}

% When standalone: set counter so \chapter{} renders as "BAB 6".
% When in main.tex: main.tex sets the counter before \subfile{bab-6}.
\ifSubfilesClassLoaded{}{\setcounter{chapter}{5}}

\chapter{Penutup}\label{sec:penutup}

Setelah artefak dibangun dan dievaluasi pada bab-bab sebelumnya, penelitian ini
sampai pada tahap penyimpulan. Bab ini menautkan kembali temuan evaluasi pada
Bab~5 dengan keempat pertanyaan penelitian yang dirumuskan di awal, menimbang
sejauh mana tujuan penelitian telah tercapai sekaligus mengakui keterbatasan
yang masih melekat, sebelum menguraikan arah pengembangan selanjutnya.

% ============================================================
\section{Kesimpulan}
% ============================================================

Penelitian ini, dalam kerangka \textit{Design Science Research} (DSR), berhasil
merancang, membangun, dan mengevaluasi sebuah \textit{pipeline}
\textit{LLM-Assisted Knowledge Graph Completion} (KGC) untuk memetakan
interkoneksi materi sains lintas-disiplin (Fisika, Kimia, dan Biologi) pada
Fase~F Kelas~XII Kurikulum Merdeka berbahasa Indonesia. Berdasarkan hasil
evaluasi struktural dan validasi pakar yang dibahas pada Bab~5, kesimpulan yang
menjawab keempat rumusan masalah penelitian adalah sebagai berikut.

\begin{enumerate}
  \item \textbf{Representasi eksplisit interkoneksi kurikulum.}
        Interkoneksi semantik antar-materi Fisika, Kimia, dan Biologi pada
        Fase~F Kelas~XII dapat direpresentasikan secara eksplisit dan dapat
        ditelusuri dalam bentuk \textit{Knowledge Graph} (KG) berbasis buku teks.
        Representasi tersebut membuat keterkaitan lintas-disiplin yang semula
        tersebar dalam struktur mata pelajaran terpisah menjadi terdokumentasi
        sebagai relasi antarkonsep. Dengan demikian, KG yang dihasilkan mampu
        menampilkan hubungan konseptual lintas-disiplin pada korpus dan
        konfigurasi penelitian ini. Kriteria keberhasilan representasi juga
        terpenuhi karena KG memuat relasi lintas-disiplin tervalidasi yang
        menghubungkan ketiga mata pelajaran, meskipun belum dimaksudkan sebagai
        peta kurikulum yang lengkap untuk seluruh fase atau jenjang.

  \item \textbf{Relasi implisit lintas-disiplin yang teridentifikasi.}
        Setelah proses LLM-Assisted \textit{Knowledge Graph Completion},
        teridentifikasi 125 relasi lintas-buku yang merepresentasikan hubungan
        implisit antar-konsep sains yang tidak tersurat dalam satu buku. Relasi tersebut didominasi
        oleh \texttt{PRASYARAT\_UNTUK} (45{,}6\%), menandakan bahwa jembatan
        lintas-disiplin pada kurikulum kelas XII lebih banyak bersifat prasyarat
        pengetahuan. Secara pola keterhubungan, Kimia berperan sebagai poros
        (\textit{hub}) yang menghubungkan Biologi (78\% relasi) dan Fisika
        (Gambar~\ref{fig:bridge-graph}). Karena tahap pencarian kandidat
        bertumpu pada kemiripan makna deskripsi, bukan pada pola konektivitas
        graf, relasi yang ditemukan terutama mencerminkan kedekatan semantik
        antar-konsep lintas-buku.

  \item \textbf{Perubahan kualitas struktural KG sebelum dan sesudah completion.}
        Kualitas struktural KG meningkat setelah proses \textit{completion},
        terutama pada kepadatan hubungan antarkonsep dan keterhubungan
        lintas-disiplin. Sebelum KGC, ketiga mata pelajaran membentuk komunitas
        yang relatif terpisah, ditandai nilai \textit{Modularity} 0{,}8906
        berdasarkan partisi komunitas hasil algoritma Louvain, yang merefleksikan
        kondisi \textit{siloed learning}. Setelah penambahan
        relasi lintas-buku, \textit{Modularity} turun menjadi 0{,}8209 dan jumlah
        komunitas terdeteksi berubah dari 19 menjadi 24; kenaikan jumlah
        komunitas ini dibaca sebagai pembagian Louvain yang lebih rinci pada graf
        pasca-KGC, bukan sebagai peningkatan keterpisahan, sementara
        \textit{Average Degree} (AD) meningkat 98{,}4\% (dari 0{,}74 menjadi
        1{,}47). Selain itu, \textit{Typed Relation Ratio} (TRR) meningkat dari
        69{,}69\% menjadi 74{,}86\%, menunjukkan bahwa lebih banyak relasi dalam
        graf memiliki tipe hubungan yang eksplisit, sedangkan
        \textit{Description Completeness} (DC) tetap 93{,}07\% karena metrik ini
        mengukur kelengkapan deskripsi \textit{node} yang tidak berubah akibat
        penambahan \textit{edge} KGC. Arah perubahan AD dan
        \textit{Modularity} tersebut memenuhi kriteria keberhasilan yang
        ditetapkan sebelum evaluasi struktural. Bukti topologis ini
        (Tabel~\ref{tab:struktural}) menunjukkan bahwa proses
        \textit{completion} memperkuat keterhubungan konseptual antar-disiplin.
        Namun, metrik graf tersebut digunakan sebagai indikator kuantitatif
        perubahan struktur, bukan sebagai bukti kebenaran semantik setiap relasi.
        Validitas relasi tetap bergantung pada hasil validasi pakar yang
        dirangkum pada poin berikutnya.

  \item \textbf{Validitas relasi dan keandalan evaluasi.}
        Relasi yang dihasilkan memiliki tingkat validitas tinggi berdasarkan
        validasi pakar, sementara keandalan evaluasinya ditopang oleh kombinasi
        metrik kuantitatif dan mekanisme \textit{LLM-as-a-judge}. Pada validasi
        intra-mapel (Fase~1) oleh enam pakar domain, sistem mencatatkan
        \textit{Fuzzy Precision} 93{,}1\%, \textit{Fuzzy Recall} 91{,}8\%, dan
        F1-Score 92{,}4\% (Tabel~\ref{tab:validasi-metrik}), dengan kesepakatan
        antarpakar Gwet's AC1 yang bervariasi dari ``baik'' hingga ``sangat
        baik'' (0{,}607 pada Kimia hingga 0{,}931 pada Biologi). Pada validasi
        pakar Fase~2 atas relasi lintas-buku, sistem mencapai validitas 99{,}1\%
        (116/117) dan ketepatan tipe 98{,}3\% (115/117), tetapi ketepatan arah
        hanya 75{,}9\% (88/116). Mekanisme \textit{LLM-as-a-judge}
        (Gemini~3.1 Pro Preview) menengahi 123 konflik rating, memilih salah satu
        label yang telah diberikan pakar pada 95{,}1\% kasus alih-alih menilai
        ulang secara independen, serta menyetujui 29 dari 44 usulan
        \textit{missing triple} (Tabel~\ref{tab:llm-judge}). Dengan demikian,
        KG yang dihasilkan valid secara relasional pada konfigurasi yang diuji:
        validitas Fase~2 melampaui ambang 90\% yang ditetapkan sebelum
        evaluasi, sementara seluruh nilai AC1 Fase~1 memenuhi ambang minimal
        0{,}6. Perlu dicatat bahwa AC1 tidak dihitung untuk Fase~2 karena fase
        tersebut menggunakan diskusi panel, bukan dua penilaian independen per
        relasi; normalisasi arah relasi masih menjadi kelemahan utama yang perlu
        diperbaiki.
\end{enumerate}

Secara keseluruhan, penelitian ini menunjukkan bahwa pendekatan LLM-Assisted KGC
merupakan metode yang layak (\textit{viable}) untuk membangun dan melengkapi graf
pengetahuan kurikulum sains lintas-disiplin berbahasa Indonesia pada korpus dan
konfigurasi yang diuji. Artefak yang
dihasilkan berupa purwarupa KG tervalidasi beserta \textit{pipeline} yang
terdokumentasi, dan menjadi salah satu sumber daya terbuka berbahasa Indonesia
yang mendokumentasikan interkoneksi materi sains lintas-disiplin pada Kurikulum
Merdeka. Artefak ini dapat dimanfaatkan oleh pengembang kurikulum dan peneliti
teknologi pendidikan, dengan catatan
penting bahwa normalisasi arah relasi masih memerlukan penyempurnaan lebih lanjut.

% ============================================================
\section{Keterbatasan Penelitian}
% ============================================================

Penelitian ini memiliki sejumlah keterbatasan yang perlu diperhatikan dalam
menafsirkan hasilnya.

\begin{enumerate}
  \item \textbf{Cakupan korpus.} Korpus terbatas pada buku teks siswa beserta
        Buku Panduan Guru pasangannya untuk Fisika, Kimia, dan Biologi Fase~F
        Kelas~XII (\textit{e-book} resmi Kemendikbudristek). Buku teks Kelas~XI
        tidak disertakan, sehingga generalisasi temuan ke fase atau jenjang lain
        belum dapat dipastikan.

  \item \textbf{Penegakan vokabulari relasi.} Tipe relasi dalam-buku diarahkan
        melalui \textit{prompt} dan tidak ditegakkan secara kaku pada level
        skema; sekitar 25\% \textit{edge} dalam-buku memuat tipe di luar 15 tipe
        inti ontologi. Tipe-tipe tersebut dikonsolidasikan pada tahap konsensus,
        namun ketiadaan penegakan keras (mis.\ melalui \texttt{enum}) merupakan
        keterbatasan reprodusibilitas dan dapat menggelembungkan keragaman tipe
        relasi.

  \item \textbf{Kesalahan arah relasi.} Sebagaimana disimpulkan di atas, sekitar
        seperempat relasi lintas-buku memiliki arah yang terbalik. Analisis
        \textit{parameter sweep} (Tabel~\ref{tab:validasi-sweep}) menunjukkan bahwa
        kesalahan ini tidak teratasi dengan pengetatan \textit{threshold},
        melainkan memerlukan perbaikan pada \textit{prompt} klasifikasi atau
        normalisasi arah pasca-proses yang belum diimplementasikan pada
        penelitian ini.

  \item \textbf{Pengaruh relasi audit berlabel \texttt{wrong}.} Tahap
        \textit{completion} dijalankan di atas graf konsensus yang masih
        mempertahankan relasi intra-mapel berlabel \texttt{wrong} sebagai
        artefak audit. Relasi tersebut tidak memengaruhi pencarian kandidat ANN
        secara langsung karena \textit{embedding} dibangun dari deskripsi
        \textit{node}; namun relasi tersebut berpotensi muncul sebagai konteks
        graf pada classifier LLM. Audit dampak pada Bab~5 menunjukkan bahwa 15
        dari 125 \textit{edge} \textit{completion} melibatkan konsep sumber
        relasi \texttt{wrong}; 13 di antaranya masuk validasi Fase~2 dan
        seluruhnya dinilai valid serta tepat tipe, meskipun 5 memiliki arah
        terbalik. Karena itu, dampaknya dibaca sebagai keterbatasan yang
        terutama berkaitan dengan normalisasi arah, bukan sebagai pembatal
        validitas relasi secara umum.

  \item \textbf{\textit{Under-extraction} relasi kuantitatif.} LLM cenderung kuat
        pada konsep deskriptif namun lemah pada relasi berbasis persamaan
        matematis; 67\% \textit{missing triples} yang diusulkan pakar Fisika
        berupa rumus yang tidak berhasil diekstrak
        (Subbab~\ref{subsubsec:error-analysis}).

  \item \textbf{Batasan inferensi pedagogis.} Sistem hanya menginferensi relasi
        \textit{single-hop} dan tidak menyusun jalur belajar multi-langkah
        lintas-disiplin maupun urutan pembelajaran yang terkomposisi.

  \item \textbf{Keterbatasan metrik dan independensi evaluasi.} Validitas relasi
        lintas-disiplin (Fase~2, 99{,}1\%) berasal dari \emph{satu} keputusan
        konsensus panel yang dinilai bersama; akibatnya \textit{Fuzzy Recall},
        F1-Score, maupun kesepakatan antarpakar (Gwet's AC1) tidak dapat dihitung
        untuk fase ini, dan angka validitas tersebut perlu dibaca dengan
        hati-hati karena tidak disertai ukuran keandalan serta rentan terhadap
        \textit{leniency} dan dominasi suara dalam diskusi panel. Selain itu,
        model ekstraksi, \textit{completion}, dan \textit{LLM-as-a-judge} berasal
        dari satu keluarga (Gemini) tanpa komparasi lintas-model, sehingga
        \textit{judge} tidak sepenuhnya independen dari sistem yang dinilainya.
        Graf juga diinstansiasi pada skala purwarupa (Neo4j Aura \textit{Free
        tier}) tanpa integrasi ke sistem pengguna akhir.

  \item \textbf{Ketiadaan metode pembanding (\textit{baseline}).} Klaim bahwa pendekatan berbasis
        LLM lebih sesuai daripada metode \textit{embedding} struktural untuk
        konsep \textit{long-tail} (Bab~2) bertumpu pada literatur dan tidak diuji
        langsung pada korpus ini: penelitian tidak membandingkan
        \textit{pipeline} terhadap baseline mana pun (misalnya \textit{embedding}
        struktural seperti TransE/RotatE, ekstraksi berbasis aturan, atau LLM
        tanpa penyaringan kandidat ANN), dan belum melakukan analisis kuantitatif
        atas distribusi \textit{long-tail}. Karena itu, klaim kelayakan bersifat
        absolut pada konfigurasi yang diuji, bukan komparatif.

  \item \textbf{Profil dan kedalaman pengalaman pakar.} Validasi dilakukan oleh
        mahasiswa aktif atau lulusan baru (angkatan 2022--2023) berlatar ilmu
        murni atau kependidikan yang linier (\ref{lamp:profil-pakar}), bukan
        pakar kurikulum atau praktisi senior. Karena penilaian ``relevansi
        pedagogis'' bergantung pada kedalaman pengalaman ini, klaim pedagogis
        sebaiknya ditafsirkan sebagai penilaian keilmuan per-relasi, bukan
        rekomendasi kurikuler yang telah tervalidasi dalam praktik pembelajaran.

  \item \textbf{Kualitas data konsep rintisan.} Sebanyak 25 \textit{node}
        \texttt{Concept} berdeskripsi kosong (tercermin pada DC 93{,}07\%) dan
        berupa istilah rintisan huruf kecil. Karena tetap berlabel
        \texttt{Concept}, konsep-konsep ini ikut di-\textit{embed} dan
        dipertimbangkan pada tahap \textit{completion} meskipun deskripsinya
        tidak memadai, sehingga dapat memunculkan relasi lintas-buku
        \texttt{BERKAITAN\_DENGAN} berkeyakinan rendah. Pembersihan atau
        pelengkapan deskripsi konsep rintisan ini belum dilakukan dan menjadi
        arah perbaikan kualitas data.
\end{enumerate}

% ============================================================
\section{Saran}
% ============================================================

Berdasarkan kesimpulan dan keterbatasan di atas, beberapa saran diajukan untuk
penelitian dan pengembangan selanjutnya.

\subsection{Bagi Peneliti dan Pengembang Teknologi Pendidikan}

\begin{enumerate}
  \item \textbf{Penanganan arah relasi.} Memprioritaskan perbaikan ketepatan arah,
        misalnya melalui penambahan contoh \textit{few-shot} yang eksplisit
        menunjukkan arah, langkah klasifikasi arah terpisah, atau normalisasi
        arah pasca-proses berbasis \textit{reasoning} semantik. Mekanisme ini
        dapat memeriksa apakah arah yang dipilih sudah sesuai dengan makna
        relasi, misalnya apakah konsep sumber benar-benar berperan sebagai
        prasyarat, aplikasi, atau pendalaman dari konsep tujuan, sehingga
        koreksi arah tidak hanya mengikuti pola dominan tertentu dalam data.

  \item \textbf{Pemisahan graf audit dan graf tervalidasi.} Mengembangkan
        \textit{validated graph view} sebelum tahap \textit{completion}, sehingga
        relasi dengan \texttt{consensus}/\texttt{status} berlabel
        \texttt{wrong} tetap tersedia untuk audit dan analisis galat, tetapi
        dikeluarkan dari konteks classifier serta mekanisme
        \textit{skip-existing-edge}. Perbandingan completion pada graf audit dan
        graf tervalidasi dapat digunakan untuk mengukur dampak relasi salah
        secara lebih langsung.

  \item \textbf{Penguatan ekstraksi relasi kuantitatif.} Menambahkan tahap
        ekstraksi khusus untuk rumus dan persamaan matematis, atau memanfaatkan
        kemampuan \textit{multimodal} model untuk membaca notasi matematis dan
        diagram, guna mengurangi \textit{under-extraction}.

  \item \textbf{Perluasan korpus dan generalisasi.} Memperluas korpus ke buku teks
        Kelas~XI, dokumen Capaian Pembelajaran (CP), serta mata pelajaran dan fase
        lain untuk menguji generalisasi \textit{pipeline}.

  \item \textbf{Inferensi multi-langkah.} Mengembangkan kemampuan penalaran
        \textit{multi-hop} untuk menyusun jalur belajar lintas-disiplin
        terkomposisi yang adaptif.

  \item \textbf{Evaluasi lintas-model.} Membandingkan beberapa keluarga LLM (mis.\
        GPT, Claude) pada tahap ekstraksi, \textit{completion}, dan
        \textit{LLM-as-a-judge} untuk menguji ketangguhan hasil, serta memperkuat
        validasi Fase~2 dengan dua pakar independen agar \textit{Fuzzy Recall} dan AC1
        dapat dihitung.
\end{enumerate}

\subsection{Bagi Pendidik dan Pengembang Kurikulum}

KG yang dihasilkan dapat dimanfaatkan sebagai referensi dalam merancang unit
pembelajaran terpadu lintas-disiplin, khususnya dengan memanfaatkan posisi Kimia
sebagai poros konseptual yang menjembatani Biologi dan Fisika. Pengembangan
lanjutan dapat membangun lapisan aplikasi bagi pengguna akhir di atas KG ini,
seperti perkakas perancangan kurikulum, sistem rekomendasi pembelajaran adaptif,
atau integrasi ke dalam platform LMS, yang berada di luar lingkup penelitian ini.

\ifSubfilesClassLoaded{\printbibliography[title={Daftar Pustaka}]}{}

\end{document}
